\chapter{项目背景与机会判断}

\section{政策环境与时代窗口}

智舆项目的诞生源于三个层面的机会交汇：AI技术的成熟、城市治理的数字化转型需求、以及当前政策环境的有力支撑。国家层面，数字政府、智慧城市、人工智能+、社会治理现代化和新质生产力等政策持续释放需求信号；省市层面，河南正加快推进数字化治理体系建设，信阳在文旅服务、基层治理、热线协同和城市品牌管理方面均存在对低成本智能化工具的现实需求。智舆不仅具备技术创新空间，也具备明确的政策承接场景。

\begin{figure}[htbp]
	\centering
	\includegraphics[width=0.88\textwidth]{chapter1-policy-system.png}
	\caption{国家-省-市三级政策支持体系图}
\end{figure}

从政策层级看，国家层面给出智慧城市、数字政府与人工智能+的总方向，河南层面给出城市全域数字化转型和人工智能产业扶持口径，信阳层面则已经形成数字政府、12345 热线和基层治理数字化的现实基础。该三级结构说明，智舆并不是脱离政策场景的抽象技术方案，而是能够与地方治理数字化升级形成直接衔接的应用产品。

\section{传统舆情治理的核心痛点}

\begin{table}[htbp]
	\centering
	\caption{传统舆情治理主要痛点}
	\begin{tabularx}{\textwidth}{L{3cm}YY}
		\toprule
		\tblhead{痛点类别} & \tblhead{具体表现}              & \tblhead{对业务造成的影响}  \\
		\midrule
		\rowcolors{2}{tablealt}{white}
		数据覆盖不足         & 平台来源分散，图片、短视频、评论和热线数据难以统一纳入 & 风险发现滞后，容易漏掉早期弱信号    \\
		分析依赖人工         & 研判人员需要手工筛选、分类、汇总与写报         & 响应速度慢，人工成本高，结论一致性不足 \\
		缺少处置建议         & 系统能报警，但难以形成面向管理者的处置方案比较     & 监测后仍然要依赖经验判断，决策链路长  \\
		下沉市场门槛高        & 成熟系统价格高、部署重、服务半径集中在一二线城市    & 地市县客户难以承担，长期存在工具缺口  \\
		空间表达不足         & 传统报表偏列表化、平面化，难以结合城市空间态势理解问题 & 指挥调度和跨部门协同效率受限      \\
		\bottomrule
	\end{tabularx}
\end{table}

从市场空间看，项目所对应的并不是一个收缩中的存量行业。根据现有背景调查资料，全球舆情监测系统市场 2024 年约为 23.15 亿美元，预计 2031 年将达到 40.50 亿美元；中国市场 2025 年预计达到 72.4 亿元，并保持较快增长。真正的机会不只在“有没有市场”，而在于现有主流产品主要覆盖高预算客户，而地市、县域、景区、高校与中小企业场景仍缺少价格可承受、部署足够轻、表达足够直观的解决方案。

\begin{figure}[htbp]
	\centering
	\includegraphics[width=0.86\textwidth]{chapter1-market-growth.png}
	\caption{中国舆情监测市场规模增长趋势图（2020-2030）}
\end{figure}

\section{信阳场景为何适合作为切入点}

信阳兼具文旅流量、基层治理、城市品牌传播和区域发展叙事等多重特征。一方面，茶旅融合、节假日客流、景区服务与消费体验容易产生公开舆情；另一方面，基层治理、热线协同、校园管理与企业服务也要求更快的风险识别和闭环处置。对智舆而言，信阳不是抽象的样本城市，而是兼具真实业务场景、试点可达性和本地资源协同优势的首批落地环境。


智舆在信阳的试点布局覆盖景区文旅、基层民生、校园治理、品牌监测四类核心场景，形成多场景协同验证的样板体系。景区文旅场景聚焦投诉热点监测与节假日舆情波动管理；基层民生场景关注高频投诉主题识别与跨渠道重复诉求处理；校园治理场景针对校园热点事件与社媒情绪变化监测；品牌监测场景服务本地企业的品牌声量与舆论风险管理需求。四类场景相互补充，既验证系统在不同治理对象中的适配能力，也为后续区域复制积累可参考的实施样板。
